% Generated draft: requires mathematical and source review.
\documentclass[11pt]{article}
\usepackage[margin=25mm]{geometry}
\usepackage{fontspec}
\setmainfont{Latin Modern Roman}
\usepackage{amsmath,amssymb,mathtools,unicode-math}
\setmathfont{Latin Modern Math}
\usepackage{enumitem,needspace,xurl,hyperref}
\setlength{\parindent}{0pt}
\setlength{\parskip}{5pt}
\setlength{\emergencystretch}{3em}
\newcommand{\N}{\mathbb{N}}
\newcommand{\Z}{\mathbb{Z}}
\newcommand{\Q}{\mathbb{Q}}
\newcommand{\R}{\mathbb{R}}
\newcommand{\C}{\mathbb{C}}
\newcommand{\F}{\mathbb{F}}
\newcommand{\A}{\mathbb{A}}
\newcommand{\PP}{\mathbb{P}}
\newcommand{\E}{\mathbb{E}}
\newcommand{\eps}{\varepsilon}
\newcommand{\dd}{\,\mathrm{d}}
\begin{document}
\section*{UnsolvedMath ID: 30001838}
\textbf{Source problem code:} catalogue entry
\subsection{Definitions and mathematical statement}
Let \(G\) be a finite group. For \(g\in G\), its conjugacy class is \(g^G=\{xgx^{-1}:x\in G\}\), and let \(\operatorname{Cl}(G)\) denote the set of conjugacy classes of \(G\). Write \(H<G\) when \(H\) is a proper subgroup of \(G\).

The supplied question is: \emph{How many conjugacy classes of a finite group \(G\) are needed to distinguish \(G\) from each of its proper subgroups?}

The record does not define ``distinguish'' or specify how the classes are selected. Under the explicit interpretation that a class distinguishes \(G\) from \(H\) when it contains no element of \(H\), a family \(\mathcal S\subseteq\operatorname{Cl}(G)\) distinguishes \(G\) from all its proper subgroups precisely when
\[
  \forall H<G\;\exists C\in\mathcal S\quad C\cap H=\varnothing.
\]
Equivalently, no proper subgroup meets every class in \(\mathcal S\). Under this interpretation, the minimum number would be
\[
 d(G)=\min\bigl\{|\mathcal S|:\mathcal S\subseteq\operatorname{Cl}(G),\quad
 \forall H<G\;\exists C\in\mathcal S\ (C\cap H=\varnothing)\bigr\},
\]
with the convention \(d(G)=+\infty\) if the defining collection is empty. Empty families are allowed, so this convention gives \(d(\{1\})=0\).

A conditional precise formulation is therefore: \emph{Determine \(d(G)\), or bounds for it, for finite groups \(G\).} This is an interpretation requiring review, not an unambiguously specified target of the supplied wording. In particular, the supplied Chebotarev-invariant references suggest that an expected number of randomly sampled classes may instead be intended. The record specifies neither a sampling distribution nor a stopping rule, and it does not determine whether a minimum, an expectation, or another quantitative bound is requested.
\subsection{Short English statement}
The question asks how many conjugacy classes distinguish a finite group from every proper subgroup. One possible meaning requires every proper subgroup to miss at least one selected class. The draft formalizes the minimum number under that interpretation, while leaving unresolved whether the intended quantity instead concerns random sampling and a Chebotarev invariant.
\subsection{Sources}
\begin{itemize}
\item Supplied local UnsolvedMath registry record, permanent ID \(30001838\), title ``Conjugacy Classes Distinguishing Finite Groups from Subgroups,'' source problem code \texttt{catalogue entry}, in the Oberwolfach Reports Open Problems collection. The record gives proposed year 2011 and identifies no proposer.
\item Supplied report citation: \emph{Explicit Methods in Number Theory (2012)}, Oberwolfach Reports, DOI 10.4171/owr/2011/35, pp. 1971--2044. \url{https://doi.org/10.4171/owr/2011/35}. The background describes Oberwolfach Reports as documenting mathematical workshops and their open-problem discussions. The year 2012 in the supplied title and the 2011 provenance are retained without reconciliation.
\item Supplied literature reference, labeled ``Kowalski--Zywina Chebotarev invariant'': \url{https://arxiv.org/abs/1008.4909}.
\item Supplied literature reference, labeled ``later conjecture and partial results'': \url{https://arxiv.org/abs/1509.05859}.
\item Literature assessment checked at \texttt{2026-08-21}: the supplied curation labels the problem partially solved and states that the Chebotarev-invariant framework and substantial finite-group results are known, while the extracted general question lacks a complete uniform answer. This is supplied dated curation derived from the OWR material, not an independently verified literature assessment; time-sensitive status and cited claims require independent verification.
\item The registry labels the statement \texttt{corrected\_verified} and reports that its mathematical content was checked against the official Oberwolfach report, with notation, context, and question form normalized without changing the source-level claim. This is a supplied verification claim; no independent source inspection is asserted here.
\end{itemize}
\end{document}
